% tgp_ppn_full.tex -- Complete PPN parameter set for TGP
% Self-contained section for \input{} into main manuscript.
% Reference standard: Will (2014), "The Confrontation between GR and Experiment"
% Synchronized with sek08c_metryka_z_substratu.tex (thm:metric-from-substrate-full).

\subsection{Complete PPN parameter set}%
\label{ssec:ppn-full}

We derive the full set of ten PPN parameters for TGP following the
standard framework of Will~(2014).  The starting point is the
effective metric (Thm.~\ref{thm:metric-from-substrate-full}):
\begin{equation}\label{eq:metric-ppn}
  ds^2 = -c_0^2\,e^{-2U/c_0^2}\,dt^2 + e^{+2U/c_0^2}\,\delta_{ij}\,dx^i dx^j,
\end{equation}
where $U$ is the Newtonian gravitational potential satisfying
$\nabla^2 U = -4\pi G_0\,\rho$.

% ------------------------------------------------------------------
\subsubsection{Step 1: Classification as a metric theory}

\begin{proposition}[TGP is a metric theory]\label{prop:metric-theory}
In TGP, matter (including non-gravitational fields) couples
\textbf{universally} to the effective metric $g^{\rm eff}_{\mu\nu}$
and not directly to the scalar field $\Phi$.  Therefore TGP is a
\textbf{metric theory} in the sense of Will~(2014, \S3.1).
\end{proposition}

\noindent\emph{Justification.}
The TGP matter action is
$S_{\rm matter} = \int \mathcal{L}_{\rm matter}(g^{\rm eff}_{\mu\nu}, \psi_{\rm matter})\,
\sqrt{-g^{\rm eff}}\,d^4x$,
where $g^{\rm eff}_{\mu\nu}$ depends on $\Phi$ but $\psi_{\rm matter}$
does not couple to $\Phi$ except through $g^{\rm eff}_{\mu\nu}$.
This is the defining property of a metric theory.

\medskip
For any metric theory, Will~(2014, \S4.2) proves:
\begin{equation}\label{eq:metric-theory-constraints}
  \boxed{
    \xi = \alpha_3 = \zeta_1 = \zeta_2 = \zeta_3 = \zeta_4 = 0.
  }
\end{equation}
This eliminates 6 of the 10 PPN parameters.
The remaining four are $(\gamma, \beta, \alpha_1, \alpha_2)$.

% ------------------------------------------------------------------
\subsubsection{Step 2: $\gamma_{\rm PPN}$ and $\beta_{\rm PPN}$}

From the exponential metric~\eqref{eq:metric-ppn}:
\begin{align}
  g_{tt} &= -c_0^2\,e^{-2U/c_0^2}
  = -c_0^2\bigl(1 - 2U/c_0^2 + 2U^2/c_0^4 - \ldots\bigr), \\
  g_{ij} &= e^{+2U/c_0^2}\,\delta_{ij}
  = \bigl(1 + 2U/c_0^2 + 2U^2/c_0^4 + \ldots\bigr)\delta_{ij}.
\end{align}
Comparing with the PPN expansion (Will 2014, eq.~4.48):
\begin{align}
  g_{tt}^{\rm PPN} &= -(1 - 2U + 2\beta\,U^2 + \ldots), \\
  g_{ij}^{\rm PPN} &= (1 + 2\gamma\,U)\,\delta_{ij},
\end{align}
we read off:
\begin{equation}\label{eq:gamma-beta}
  \boxed{\gamma_{\rm PPN} = 1, \qquad \beta_{\rm PPN} = 1}
  \qquad \text{(exact to all orders in $U$)}.
\end{equation}

\noindent\textbf{Note.}
The power-law identification $\psi = 1 + 2U$ gives $\beta = 2$
(see Thm.~\ref{thm:metric-from-substrate-full}, Krok~2).
The exponential form is selected by the PPN consistency
condition $\beta = 1$ (assumption~(iii) of the theorem).
Both forms agree at $O(U)$ and differ only at $O(U^2)$.

% ------------------------------------------------------------------
\subsubsection{Step 3: Preferred-frame parameters $\alpha_1$, $\alpha_2$}

The preferred-frame parameters measure effects proportional to
the velocity $\mathbf{w}$ of the PPN coordinate system relative
to the ``mean rest frame of the universe'' (Will 2014, \S4.2).
They enter the metric at $O(U \cdot w^2/c^2)$ and $O(U \cdot w/c)$
respectively.

\paragraph{Semi-conservative theories.}
For any \emph{semi-conservative} metric theory ($\zeta_i = 0$,
which TGP satisfies by Prop.~\ref{prop:metric-theory}), the
preferred-frame parameters are determined by the cosmological
time evolution of the scalar field:
\begin{equation}\label{eq:alpha-from-phi-dot}
  \alpha_1 = -\frac{4\,\dot\phi_0}{H_0\,\phi_0}\,\frac{\partial\gamma}{\partial\phi_0},
  \qquad
  \alpha_2 = -\frac{\dot\phi_0}{H_0\,\phi_0}\,\frac{\partial(\gamma - \beta - 1)}{\partial\phi_0},
\end{equation}
where $\phi_0$ is the cosmological background value of the scalar
field and $\dot\phi_0 \sim H_0\,\phi_0$ its rate of change
(Damour \& Esposito-Far\`ese 1992; Nordtvedt 1970).

\paragraph{Application to TGP.}
In TGP, $\gamma = 1$ and $\beta = 1$ are \textbf{exact} ---
they do not depend on $\Phi$ or any coupling parameter.
Therefore:
\begin{equation}
  \frac{\partial\gamma}{\partial\phi_0} = 0,
  \qquad
  \frac{\partial(\gamma - \beta - 1)}{\partial\phi_0} = 0,
\end{equation}
and consequently:
\begin{equation}\label{eq:alpha12}
  \boxed{\alpha_1 = 0, \qquad \alpha_2 = 0.}
\end{equation}

\paragraph{Physical interpretation.}
The vanishing of $\alpha_1$ and $\alpha_2$ reflects the fact that
TGP's exponential metric~\eqref{eq:metric-ppn} is structurally
identical to GR at \emph{every} post-Newtonian order.
The scalar field $\Phi$ determines the metric conformally,
but the conformal factor is fixed by the local matter distribution
(Poisson equation) --- there is no free propagating scalar degree
of freedom in the PPN regime that could ``remember'' the cosmic
rest frame.

\paragraph{Consistency check: Cassini bound.}
The Cassini spacecraft measured $\gamma - 1 = (2.1 \pm 2.3) \times 10^{-5}$
(Bertotti et al.\ 2003).
TGP: $\gamma - 1 = 0$ exactly.
\textsc{Pass.}

Lunar Laser Ranging constrains $|4\beta - \gamma - 3| < 5.3 \times 10^{-4}$
(Nordtvedt effect; Williams et al.\ 2009).
TGP: $4\beta - \gamma - 3 = 4 - 1 - 3 = 0$ exactly.
\textsc{Pass.}

Binary pulsar timing constrains $|\hat\alpha_1| < 1.3 \times 10^{-4}$,
$|\hat\alpha_2| < 1.8 \times 10^{-5}$ (Damour \& Esposito-Far\`ese 1992;
Shao et al.\ 2013).
TGP: $\alpha_1 = \alpha_2 = 0$.
\textsc{Pass.}

% ------------------------------------------------------------------
\subsubsection{Summary: complete PPN table}

\begin{center}
\begin{tabular}{@{}llll@{}}
\toprule
\textbf{Parameter} & \textbf{GR value} & \textbf{TGP value} & \textbf{Derivation} \\
\midrule
$\gamma$   & 1 & 1 & Exponential metric, eq.~\eqref{eq:gamma-beta} \\
$\beta$    & 1 & 1 & Exponential metric, eq.~\eqref{eq:gamma-beta} \\
$\xi$      & 0 & 0 & Metric theory, eq.~\eqref{eq:metric-theory-constraints} \\
$\alpha_1$ & 0 & 0 & $\partial\gamma/\partial\phi_0 = 0$, eq.~\eqref{eq:alpha12} \\
$\alpha_2$ & 0 & 0 & $\partial(\gamma{-}\beta{-}1)/\partial\phi_0 = 0$, eq.~\eqref{eq:alpha12} \\
$\alpha_3$ & 0 & 0 & Metric theory, eq.~\eqref{eq:metric-theory-constraints} \\
$\zeta_1$  & 0 & 0 & Metric theory, eq.~\eqref{eq:metric-theory-constraints} \\
$\zeta_2$  & 0 & 0 & Metric theory, eq.~\eqref{eq:metric-theory-constraints} \\
$\zeta_3$  & 0 & 0 & Metric theory, eq.~\eqref{eq:metric-theory-constraints} \\
$\zeta_4$  & 0 & 0 & Metric theory, eq.~\eqref{eq:metric-theory-constraints} \\
\bottomrule
\end{tabular}
\end{center}

\medskip
\noindent\textbf{Result:}
TGP is \textbf{PPN-indistinguishable from GR} at all post-Newtonian
orders probed by current experiments.  All 10 PPN parameters take
their GR values \emph{exactly} (not as a limit of some coupling
constant).  This is a structural consequence of: (i)~the metric
theory property, (ii)~the exponential conformal structure,
(iii)~the fact that $\gamma$ and $\beta$ are field-independent.

\paragraph{Where TGP departs from GR.}
The PPN identity does \emph{not} mean TGP $=$ GR.  Deviations appear:
\begin{enumerate}[nosep]
  \item \textbf{Cosmological scales:} Modified Friedmann equation
        with $G(\Phi)$ and $\Lambda_{\rm eff} = \gamma/12$.
  \item \textbf{Strong-field regime:} Frozen interior (no singularity)
        instead of Schwarzschild/Kerr singularity.
  \item \textbf{Additional polarization:} Breathing mode GW (scalar,
        screened by $m_{\rm sp} \sim H_0$).
  \item \textbf{Three-body forces:} $V_3$ from $\Phi^4$ self-interaction
        (no analog in GR).
  \item \textbf{Particle physics:} Lepton mass hierarchy from soliton
        ODE ($r_{21} = 206.77$).
\end{enumerate}

\paragraph{Caveat: Nordtvedt effect and strong-field corrections.}
The analysis above is strictly valid in the \emph{weak-field}
PPN regime ($U/c^2 \ll 1$).  For neutron stars, where
$U/c^2 \sim 0.2$, strong-field corrections (``sensitivities'')
could modify the effective PPN parameters.  In standard
scalar-tensor theories, the Nordtvedt effect vanishes when
$\gamma = \beta = 1$, which TGP satisfies.
A fully rigorous strong-field analysis (following Damour \&
Esposito-Far\`ese 1993) requires solving the TGP field
equation inside a neutron star --- this is left for future work.
